\section{Law of Large Number}

\subsection{Independence}

\begin{theorem}
    [$\pi$-$\lambda$] For some families of subsets $\mathcal{A},\mathcal{L}$, we say that $\mathcal{A}$ is a $\pi$-system if it is closed under intersection. We say $\mathcal{L}$ is a $\lambda$-system if
    \begin{enumerate}
        \item $\Omega \in \mathcal{L}$
        \item If $A,B \in \mathcal{L}$ and $A\subset B$, then $B-A \in \mathcal{L}$.
        \item If $A_n \in \mathcal{L}$ and $A_n \uparrow A$, then $A\in\mathcal{L}$.
    \end{enumerate}
    Then if $\mathcal{A}$ is a $\pi$-system and $\mathcal{L}$ is a $\lambda$-system containing $\mathcal{A}$, then $\sigma(\mathcal{A}) \subset \mathcal{L}$.
\end{theorem}

\begin{theorem}
\label{thm: independence_of_generated_sigma}
    Suppose $\mathcal{A}_1, \mathcal{A}_2, \cdots , \mathcal{A}_n$ are independent and each $\mathcal{A}_i$
is a $\pi$-system. Then $\sigma(\mathcal{A}_1), \sigma(\mathcal{A}_2), \cdots , \sigma(\mathcal{A}_n)$ are independent.
\end{theorem}
\begin{proof}
    Consider
    \[
    \mathcal{M} = \left\{A \subset \Omega| P\left(A\prod_{i\in I} A_i\right) = P(A)\prod_{i=I}P(A_i)\text{ for any }I\subset\{2,3,\cdots,n\},A_i \in \mathcal{A}_i \right\},
    \]
    then $A_1 \subset \mathcal{M}$ and $\Omega \in \mathcal{M}$ in an obvious way. Assume that $A,B \in \mathcal{M}$, then
    \[
    \begin{aligned}
    P\left((B-A)\prod_{i\in I} A_i\right) &= P\left(B\prod_{i\in I} A_i - A\prod_{i\in I} A_i\right) \\ &= \left(P(B)-P(A)\right)\prod_{i\in I} P(A_i) \\ &= P(B-A)\prod_{i\in I} P(A_i).
    \end{aligned}
    \]
    For $A_n \in \mathcal{M}$ with $A_n \uparrow A$, we know
    \[
    \begin{aligned}
    P\left(A\prod_{i\in I} A_i\right) &= P\left(\bigcup_{n\geq 1} \left(A_n \prod_{i\in I} A_i\right) \right) \\ &= \limsup \left(P(A_n)\prod_{i\in I}P\left(A_i\right)\right) \\ &= P(A)\prod_{i\in I} P(A_i).
    \end{aligned}
    \]
    and hence $\mathcal{M}$ is a $\lambda$-system. Therefore, $\sigma(\mathcal{A}_1)\subset \mathcal{M}$ which is equivalent with that \[\sigma(\mathcal{A}_1), \mathcal{A}_2,\cdots,\mathcal{A}_n \text{ are independent}.\]
    For the remaining, we may prove in a similar way by replacing $\mathcal{A}_1$ by $\sigma(\mathcal{A}_1)$.
\end{proof}

\begin{proposition}
    In order for $X_1, \cdots , X_n$ to be independent, it is sufficient that for all $x_1, \cdots , x_n \in (−\infty,\infty]$, we have
\[P(X_1 \leq x_1, \cdots , X_n \leq x_n) = \prod_{i=1}^n P(X_i \leq x_i)\]
\end{proposition}
\begin{proof}
    Let $\mathcal{A}_i = \{\{X_i \leq a\}, a\in\mathbb{R}\}$, and then we know $\mathcal{A_i}$ are independent and by the Theorem \ref{thm: independence_of_generated_sigma}, $\sigma(\mathcal{A}_i)$ are independent and since $\sigma(\mathcal{A}_i) = \sigma(X_i)$ and we are done.
\end{proof}

\begin{theorem}
\label{thm: independent-algebras-matrix}
    Suppose $\mathcal{F}_{i,j} , 1 \leq i \leq n, 1 \leq j \leq m(i)$ are independent $\sigma$-algebras and let $\mathcal{G}_i = \sigma\left(\bigcup_j \mathcal{F}_{i,j}\right)$. Then $\mathcal{G}_1, \cdots , \mathcal{G}_n$ are independent.
\end{theorem}
\begin{proof}
    Let $\mathcal{A}_i = \{\cap_{j} A_{i,j}, A_{i,j}\in\mathcal{F}_{i,j}\}$ and then we know $\cup_j \mathcal{F}_{i,j} \subset \mathcal{A}_i$ and $\mathcal{A}_i$ is a $\pi$-system. Also, we would like to see that $\mathcal{A_i}$ are independent. Notice that $\Omega \in \mathcal{A_i}$ and
    \[
    \begin{aligned}
        P\left(\prod_i A_i\right) &= P\left(\prod_i\left(\cap_{j}A_{i,j}\right)\right) \\
        &= P\left(\prod_{i,j} A_{i,j}\right) \\
        &= \prod_{i}\left(\prod_j P(A_{i,j})\right) \\
        &= \prod_i P(A_i)
    \end{aligned}
    \]
    where $A_{i,j}$ satisfy that $A_i = \cap_j A_{i,j}$. Therefore, we know that $\mathcal{A}_i$ are independent and we are done by noticing that $\mathcal{G}_i \subset \sigma(\mathcal{A}_i)$ and the Theorem \ref{thm: independence_of_generated_sigma}.
\end{proof}

\begin{theorem}
\label{thm: independent-variable-matrix}
    If for $1 \leq i \leq n, 1 \leq j \leq m(i), X_{i,j}$ are independent and $f_i: \R^{m(i)} \to \R$ are measurable then $f_i(X_{i,1}, \cdots , X_{i,m(i)})$ are
independent.
\end{theorem}
\begin{proof}
    Notice that
    \[
    \sigma(f_i(X_{i,1},\cdots,X_{i,m(i)})) \subset \sigma(X_{i,1},\cdots,X_{i,j})
    \]
    and we may use the Theorem \ref{thm: independent-algebras-matrix} with $\mathcal{F}_{i,j} = \sigma(X_{i,j})$.
\end{proof}

\subsection{Distribution and Expectation for Independent Variables}

\begin{lemma}
    Suppose $X_1, \cdots , X_n$ are independent random variables
and $X_i$ has distribution $\mu_i$, then $(X_1, \cdots , X_n)$ has distribution $\mu_1 \times \cdots \times \mu_n$.
\end{lemma}

\begin{theorem}
\label{thm: expectation-of-idv}
    Suppose $X$ and $Y$ are independent and have distributions $\mu$ and $\nu$. If $h : \R^2 \to \R$ is a measurable function with $h \geq 0$ or $\mathbb{E}|h(X, Y )| < \infty$ then
    \[\mathbb{E}h(X, Y ) = \int\int h(x, y) \mu(dx) \nu(dy).\]
    In particular, if $h(x, y) = f(x)g(y)$ where $f, g : \R \to \R$ are measurable
functions with $f, g \geq 0 $ or $\mathbb{E}|f(X)|$ and $\mathbb{E}|g(Y )| < \infty$ then
\[Ef(X)g(Y ) = Ef(X) \cdot Eg(Y ).\]
\end{theorem}

\begin{proposition}
    If $X_1, \cdots , X_n$ are independent and have (a) $X_i \geq 0$ for all $i$, or (b) $\mathbb{E}|X_i| < \infty$ for all $i$ then
    \[\mathbb{E}\left[\prod_{i=1}^n X_i\right] = \prod_{i=1}^n \mathbb{E}X_i\]
\end{proposition}
\begin{proof}
    We may use the Theorem \ref{thm: independent-variable-matrix} to see that $X:= X_1$ and $Y:= X_2\cdots X_n$ are independent. By the Theorem \ref{thm: expectation-of-idv}, we know that
    \[
    \mathbb{E}XY = \mathbb{E}X\mathbb{E}Y
    \]
    and we may prove the conclusion inductively.
\end{proof}

\begin{lemma}
    If $X$ and $Y$ are independent, $F(x) = P(X \leq x)$, and $G(y) = P(Y \leq y)$, then
    \[P(X + Y \leq z) = \int F(z - y) dG(y)\]
    The integral on the right-hand side is called the convolution of $F$ and $G$ and is denoted as $F*G(z)$. 
\end{lemma}

\begin{theorem}
\label{thm: independent-density-form}
    Suppose that $X$ with density $f$ and $Y$ with distribution function $G$ are independent. Then $X + Y$ has density \[h(x) = \int f(x - y) dG(y).\] When $Y$ has density $g$, the last formula can be written as \[h(x) = \int f(x - y) g(y) dy.\]
\end{theorem}

\begin{definition}
    [Gamma and Normal Distribution] The gamma density with parameters $\alpha$ and $\lambda$ is given by
    \[
    f(x) = \begin{cases}
        \lambda^{\alpha}x^{\alpha-1}e^{-\lambda x}/ \Gamma(\alpha), \quad &x\geq 0\\
        0,&x<0
    \end{cases},
    \]
    where $\Gamma(\alpha) = \int_0^{\infty} x^{\alpha-1}e^{-x}dx$.\par
    The normal density with mean $\mu$ and variance $a$ is
    \[
    \dfrac{1}{\sqrt{2\pi a}}\exp\left({-\dfrac{(x-\mu)^2}{2a}}\right).
    \]
\end{definition}

\begin{proposition}
    If $X = gamma(\alpha, \lambda)$ and $Y = gamma(\beta, \lambda)$ are independent then $X + Y$ is $gamma(\alpha + \beta, \lambda)$. Consequently if $X_1,\cdots, X_n$ are independent $exponential(\lambda)$ r.v.’s, then $X_1 +\cdots+X_n$, has a $gamma(n, \lambda)$ distribution.
\end{proposition}
\begin{proof}
    We know the density formula can be given by
    \[
    \begin{aligned}
        h(x) &= \int^x_0 \left(\lambda^{\alpha}(x-y)^{\alpha-1}e^{-\lambda (x-y)}/ \Gamma(\alpha)\right)\left(\lambda^{\beta}y^{\beta-1}e^{-\lambda y}/ \Gamma(\beta)\right) dy \\
        &= \lambda^{\alpha+\beta}x^{\alpha+\beta - 1}e^{-\lambda x}\int_0^xx^{1-\alpha-\beta}(x-y)^{\alpha-1}y^{\beta-1}/(\Gamma(\alpha)\Gamma(\beta)) dy \\
        &=\lambda^{\alpha+\beta}x^{\alpha+\beta - 1}e^{-\lambda x}\int_0^1(1-u)^{\alpha-1}u^{\beta-1}du/(\Gamma(\alpha)\Gamma(\beta))
    \end{aligned}
    \]
    It remains to show that
    \[
    \int_0^1(1-u)^{\alpha-1}u^{\beta-1}du = \Gamma(\alpha)\Gamma(\beta)/\Gamma(\alpha+\beta)
    \]
    which can be shown by noticing
    \[
    \begin{aligned}
    \Gamma(\alpha+\beta)\int_0^1(1-u)^{\alpha-1}u^{\beta-1}du &= \int_0^{\infty}\int_0^1 v^{\alpha+\beta-1}e^{-v}(1-u)^{\alpha-1}u^{\beta-1} dvdu \\
    &=\int_0^{\infty}\int_0^1 e^{-v}(v-vu)^{\alpha-1}(vu)^{\beta - 1}dv(vdu) \\
    &= \int_0^{\infty}\int_{0}^{\infty}e^{-p+q}p^{\alpha-1}q^{\beta-1} (d(p+q)(p+q)d(q/(p+q)))\\
    &=
    \int_0^{\infty}\int_{0}^{\infty}e^{-p+q}p^{\alpha-1}q^{\beta-1}|1/(p+q)|(p+q)dpdq
    \end{aligned}
    \]
    and we are done.\par
    The remaining is obvious by notice that
    \[
    gamma(1,\lambda) = exponential(\lambda).
    \]
\end{proof}

\begin{proposition}
    If $X = normal(\mu, a)$ and $Y = normal(\nu, b)$ are independent then $X + Y = normal(\mu + \nu, a + b)$.
\end{proposition}